% Generated draft: requires mathematical and source review.
\documentclass[11pt]{article}
\usepackage[margin=25mm]{geometry}
\usepackage{fontspec}
\setmainfont{Latin Modern Roman}
\usepackage{amsmath,amssymb,mathtools,unicode-math}
\setmathfont{Latin Modern Math}
\usepackage{enumitem,needspace,xurl,hyperref}
\setlength{\parindent}{0pt}
\setlength{\parskip}{5pt}
\setlength{\emergencystretch}{3em}
\newcommand{\N}{\mathbb{N}}
\newcommand{\Z}{\mathbb{Z}}
\newcommand{\Q}{\mathbb{Q}}
\newcommand{\R}{\mathbb{R}}
\newcommand{\C}{\mathbb{C}}
\newcommand{\F}{\mathbb{F}}
\newcommand{\A}{\mathbb{A}}
\newcommand{\PP}{\mathbb{P}}
\newcommand{\E}{\mathbb{E}}
\newcommand{\eps}{\varepsilon}
\newcommand{\dd}{\,\mathrm{d}}
\begin{document}
\section*{UnsolvedMath ID: 20001640}
\textbf{Source problem code:} catalogue entry
\subsection{Definitions and mathematical statement}
The canonical statement is genuinely truncated: ``Problem 1.2. Do various candidate complexes satisfy the conditions of''. The supplied research summary says the official PDF completes the reference as Problem 1.1. Neither that problem's conditions nor the complete candidate list is included here. Therefore the main complex-comparison question cannot be recovered as a fully quantified theorem from this record, and no missing link or dynamical predicates are invented.

The aid mentions free-factor, free-splitting and cyclic-splitting complexes for $\operatorname{Out}(F_N)$, and reports an audit of their hyperbolicity and two dynamical requirements, while keeping all-link requirements unresolved. Their exact vertex/edge conventions and the two requirements are absent. This is an attributed status description, not a claim that they satisfy every omitted condition of Problem 1.1.

The supplied mathematical obstruction concerns the edge-group interpretation of a further candidate labeled $(2')$. Let $F_N$ be nonabelian free, $N\geq2$, and interpret a splitting as a Bass--Serre graph-of-groups decomposition with injective edge maps. Every nontrivial abelian subgroup of $F_N$ is cyclic. Consequently no edge group isomorphic to $\Z^r$ for $r\geq2$ can occur. Thus, if $\mathscr S_k^{\mathrm{ab}}$ is the class of splittings with free abelian edge groups of rank at most $k$, the aid states
\[
 \mathscr S_k^{\mathrm{ab}}=\mathscr S_1^{\mathrm{ab}}
 \quad\text{for all }k\geq1,
\]
and this is exactly the trivial-or-cyclic edge-group class. Allowing higher abelian rank therefore does not create a genuine filtration.

If instead the intended condition is \emph{arbitrary} edge subgroups of rank less than $k$, then $k=3$ is the first threshold allowing nonabelian edge groups. The inequality here is strict, unlike the preceding at-most-$k$ abelian convention. For every $N\geq4$, put $H=F(a,b)$ and use the displayed common-factor inclusions in
\[
 \bigl(H*\langle x\rangle*F_{N-4}\bigr)*_{H}
       \bigl(H*\langle y\rangle\bigr).
\]
The aid reports that this reduced one-edge amalgam is $F_N$ and realizes a rank-two free edge group. Both edge inclusions are proper; $F_0$ is trivial when $N=4$. This supplies an unrestricted-edge example, not an abelian rank-two splitting, and no nonexistence assertion for smaller ambient ranks is added.

These edge-rank distinctions clarify one candidate formulation but do not recover the missing source conditions or settle every complex's geometric properties.
\subsection{Short English statement}
The main question is cut off before its reference to Problem 1.1, whose conditions are absent. The supplied partial result clarifies one ambiguity: free groups have no higher-rank abelian edge subgroups, so increasing an abelian edge-rank bound stops at cyclic splittings. Allowing arbitrary edge groups is different and first permits nonabelian rank two at the stated threshold.
\subsection{Sources}
\begin{itemize}
\item Local UnsolvedMath registry, permanent ID 20001640, exact source code catalogue entry; canonical \texttt{aim-geometric-group-theory-notes.json}, \texttt{notes[148]}, source item 1.2, tag ``problem'', workshop \emph{The geometry of the outer automorphism group of a free group}; \url{https://aimath.org/WWN/outerauto/outerauto.pdf}.
\item The summary attributes completion of the missing reference to the official PDF, but the verification only introduces the truncated canonical record and no completed Problem 1.1 or candidate list is provided. No proposer or independent bibliographic references for the complex audit are supplied; original remarks and literature field are empty. This draft did not visit the PDF.
\item Supplied metadata: PARTIAL-PROGRESS, registry status ``partially\_solved'', statement status ``reconstructed\_unverified''; attempt 1, result ``partial\_result'', problem status at run ``partially\_solved'', recorded rigor ``proved'', novelty confidence ``low'', candidate type ``obstruction theorem''. The machine-generated aid is not peer reviewed and explicitly requires independent verification before citation.
\end{itemize}
\end{document}
